\section{The measurable occupation resonance group}
\label{sec:v39-resonance-group}
The occupation phase cannot be identified with a constant phase per
collision. We determine the possible structure of its measurable
obstructions before making any assertion about long powers on the full
occupation torus. This distinction also separates an arithmetic
obstruction from an estimate for the signed inverse.

All assertions in this section concern the original circular family
and its actual section. Write $c=c_*$ and use the uncentered phase
\begin{equation}\label{eq:v39-phase-equation}
 q\circ T_R=e^{i(u\cdot\kappa_R+b\tau_R+v\eta_R+s)}q,
 \qquad |q|=1\quad\hbox{almost everywhere}.
\end{equation}
Here $u\in\T^2$, $v,s\in\T$, $b\in\R$, and each torus has period
$2\pi$. A choice of representatives in the exponent is immaterial.
Let $\mathcal A_R$ be the set of $(u,b,v,s)$ for which
\eqref{eq:v39-phase-equation} has a measurable solution. Let $\mathcal G_R$
be its projection onto $(u,b,v)\in\T^2\times\R\times\T$.
These are groups of physical measurable phases, not spectra on a
postulated induced Banach space.

\subsection{A uniformly isolated group}
\begin{lemma}[Exclusion near the joint origin for measurable phases]
\label{lem:v39-local-phase-exclusion}
There is $\delta_*>0$, independent of $R$, such that
\eqref{eq:v39-phase-equation} has no solution when
\[
 0<\operatorname{dist}((u,b,v),0)<\delta_*.
\]
When $(u,b,v)=0$, it has a solution only for $s=0$; that solution is
constant almost everywhere. No regularity of $q$ is assumed.
\end{lemma}
\begin{proof}
In a coordinate neighborhood of zero put $\Theta=(u,b,v)$ and
$X_R=(\kappa_{R,1},\kappa_{R,2},\tau_R-\bar\tau_R,\eta_R-c)$.
The positive covariance and the uniform cubic expansion in
Proposition~\ref{prop:v37-joint-small-frequency} give, after reducing a
uniform radius,
\[
 |\lambda_R(\Theta)|\le e^{-a|\Theta|^2}<1
 \qquad(0<|\Theta|<\delta_*).
\]
The complementary powers decay exponentially on the same neighborhood.
Thus the exact collision pairing with any two fixed smooth endpoint
functions tends to zero at each such nonzero frequency. This assertion
uses the exact occupation multiplier, not smoothing of $X_R$.

Suppose a measurable $q$ solves \eqref{eq:v39-phase-equation}. Iteration
gives the exact identity
\begin{equation}\label{eq:v39-unit-pairing}
 \int q(x)\overline{q(T_R^m x)}
       e^{i\Theta\cdot S_{m,R}X_R(x)}\,d\nu(x)
       =e^{-im(s+b\bar\tau_R+vc)}.
\end{equation}
Choose smooth complex $q_\epsilon$, bounded by one, with
$\|q-q_\epsilon\|_1<\epsilon$. Such approximations follow by
convolution in collision charts and truncation or by a positive
smoothing of a bounded representative. Invariance and the bounded
moduli show that replacing both occurrences of $q$ changes the
integral by at most $2\epsilon$, for every $m$. With $\epsilon$ fixed,
the smooth pairing tends to zero by the preceding spectral estimate.
The right side of \eqref{eq:v39-unit-pairing} has modulus one, so
$1\le2\epsilon$, a contradiction. The approximating norm need not be
uniform in $q$ or $\epsilon$: the limit in $m$ is taken first.

At $\Theta=0$, \eqref{eq:v39-phase-equation} is the ordinary circle
eigenfunction equation. Mixing of $T_R$ excludes $s\ne0$, and
ergodicity makes an invariant $q$ constant. The mixing input is the
collision splitting, extended from smooth tests by $L^2$ approximation.
\end{proof}

\begin{theorem}[Structure and uniform separation of occupation resonances]
\label{thm:v39-resonance-structure}
For each $R$, the phase $s$ above a point of $\mathcal G_R$ is unique,
and its transfer function $q$ is unique up to a constant of modulus
one. The group $\mathcal G_R$ is closed and uniformly discrete. There
is a constant $C_*$, independent of $R$ and $B\ge0$, such that
\begin{equation}\label{eq:v39-resonance-count}
 \#\{(u,b,v)\in\mathcal G_R:|b|\le B\}\le C_*(1+B).
\end{equation}
Projection from $\mathcal G_R$ to its occupation coordinate $v$ is
injective. In particular its zero-roof subgroup is finite cyclic, of
some order $d_R\le C_*$; its occupation coordinates are exactly the
$d_R$th roots of unity in angular coordinates. Every $u$ and $s$
in that subgroup belongs to $(2\pi/d_R)\Z$ modulo $2\pi$.

There are precisely two possible forms. Either $\mathcal G_R$ is this
finite cyclic group, or its roof projection is $\beta_R\Z$ for some
$\beta_R>0$ and
\begin{equation}\label{eq:v39-resonance-structure}
 \mathcal G_R=\{n\gamma_R+j\gamma_R^{\rm tor}:
                 n\in\Z,\ 0\le j<d_R\}.
\end{equation}
Here $\gamma_R$ has roof coordinate $\beta_R$ and
$\gamma_R^{\rm tor}$ generates the zero-roof subgroup. The occupation
coordinate of $\gamma_R$, divided by $2\pi$, is irrational. The positive
roof spacing, when present, is bounded below uniformly in $R$.
No assertion that either type of nonzero resonance actually exists is
included in this classification.
\end{theorem}
\begin{proof}
Products and quotients of circle-valued solutions give addition and
subtraction of phases. If two solutions have the same $(u,b,v)$,
their quotient is a collision eigenfunction with constant phase the
difference of the two $s$'s. Mixing gives equality of those phases;
ergodicity gives constancy of the quotient. Consequently
$\mathcal A_R$ is the graph of a character $s_R$ on $\mathcal G_R$.

The difference of two distinct elements of $\mathcal G_R$ is a nonzero
element. Lemma~\ref{lem:v39-local-phase-exclusion} shows that its
distance from zero is at least $\delta_*$. Thus distinct elements are
separated by a common positive distance, without any compactness
assumption on their transfer functions. A convergent sequence in the
group is eventually constant, proving closedness. Packing disjoint
balls of radius $\delta_*/3$ into
$\T^3\times[-B-\delta_*,B+\delta_*]$ gives
\eqref{eq:v39-resonance-count}; reduce $\delta_*$ below one if needed.

If $v=0$, the phase equation is precisely the complete physical
collision phase equation. Theorem~\ref{thm:complete-physical-phase}
gives $u=0$, $b=0$, $s=0$. The kernel of projection to $v$ is therefore
trivial. The same argument applied to a difference proves injectivity.
The zero-roof subgroup is finite by \eqref{eq:v39-resonance-count},
and injects into the circle. Every finite subgroup of the circle is
cyclic and consists of all $d$th roots for its order $d$. Multiplying
any of its elements by $d$ gives zero in the full group, including
its unique $s$ coordinate. This proves the rational-coordinate claim.

Let $H_R$ be the roof projection. A bounded interval contains only
finitely many of its points, by \eqref{eq:v39-resonance-count}. Hence
this additive subgroup of $\R$ is either zero or $\beta_R\Z$ for its
least positive element. To see the latter assertion directly, choose
the least positive element in any nonempty finite set of positive
roof coordinates below a fixed bound. A smaller positive element
anywhere would be in that set. Division with remainder by this minimum
shows that every roof coordinate is its integer multiple.
Choose $\gamma_R$ over $\beta_R$. Subtracting $n\gamma_R$ from any
element with roof $n\beta_R$ leaves the finite kernel, proving
\eqref{eq:v39-resonance-structure} and uniqueness of its expression.
If the occupation coordinate of $\gamma_R$ were a rational multiple
of $2\pi$, a positive multiple of $\gamma_R$ would have zero occupation
coordinate and nonzero roof coordinate. This contradicts injectivity.
Finally, if $0<\beta_R\le1$, the elements
$0,\gamma_R,\ldots,\lfloor1/\beta_R\rfloor\gamma_R$ lie in the
unit roof slab. Their number is at most $2C_*$. This gives a positive
uniform lower bound for $\beta_R$; the case $\beta_R>1$ is immediate.
\end{proof}

\subsection{The exact arithmetic seen by an actual return}
\begin{proposition}[All-frequency induction of the phase equation]
\label{prop:v39-induced-arithmetic}
A phase $(u,b,v,s)$ belongs to $\mathcal A_R$ if and only if there is
a measurable circle-valued $q_*$ on $Y_R^*$ satisfying
\begin{equation}\label{eq:v39-induced-phase}
 q_*\circ F_R^*
   =e^{i(u\cdot\kappa_R^*+b\tau_R^*+s r_R^*+v)}q_*.
\end{equation}
This equivalence holds for every frequency, not only near zero, and
uses the actual return tower. In particular ordinary finite lattice
cover mixing does not by itself remove a nonzero $v$: on the section
$v$ is the return spectral phase.
\end{proposition}
\begin{proof}
Discard the invariant null set on which an iterate fails and telescope
\eqref{eq:v39-phase-equation} to the true first return. The occupation
sum on $[0,r_R^*)$ is exactly one. Restriction gives
\eqref{eq:v39-induced-phase}.

Conversely, write almost every collision state uniquely as $T_R^j y$,
where $y\in Y_R^*$ and $0\le j<r_R^*(y)$. With
$\psi_R=u\cdot\kappa_R+b\tau_R+s$, define
\[
 q(T_R^jy)=e^{iS_{j,R}\psi_R(y)+iv\one_{\{j\ge1\}}}q_*(y).
\]
The genuine tower is a measurable partition and this defines a
measurable circle-valued function. Its one-step equation is exact at
every nontop level. At the top it is exactly
\eqref{eq:v39-induced-phase}. This includes $r_R^*=1$, when the
initial level is already the top. No induced mixing is used.
\end{proof}

\paragraph{What is and is not removed.}
The theorem replaces an unspecified continuum of possible measurable
occupation obstructions by a uniformly separated group with finitely
many points in each bounded roof band. It also identifies their exact
return equation. It does not assert an anisotropic essential spectral
radius below one on the full occupation torus. Nor does it assert that
the finite collection on a fixed band is empty. Absence of a measurable
phase and a power estimate for an operator are different assertions.
The signed remainder in Proposition~\ref{prop:v38-exact-index-remainder}
retains its original normalization and is not set to zero by this
classification.
